\documentclass[10pt,a4paper]{amsart}
\usepackage[margin=19mm]{geometry}
\usepackage{amsmath,amssymb,mathtools}
\usepackage[scr=rsfs,cal=euler]{mathalfa}
\usepackage{xfrac}
\usepackage[unicode,hidelinks,bookmarks=false]{hyperref}
\newcommand{\ie}{i.e.\ }
\newcommand{\cf}{cf.\ }
\newcommand{\resp}{respectively\ }
\newcommand{\Subsection}{\subsection}
\providecommand{\nobreakdash}{\nobreak\hbox{-}}
\setlength{\emergencystretch}{2em}
\multlinegap0pt
\usepackage{amssymb}
\usepackage{mathtools}
\usepackage[shortlabels]{enumitem}
\usepackage{xcolor}
\newtheorem{lemma}{Lemma}
\DeclareMathOperator{\Span}{span}
\newcommand{\dd}{\,\mathrm d}
\newcommand{\norm}[1]{\left\lVert#1\right\rVert}
\title{Compactness of Toeplitz Operators on the Bergman Space}
\author{Guangfu Cao, Li He, Shuqing Zhang}
\date{Source: arXiv:2609.24614v3 (CC BY 4.0)}
\begin{document}
\maketitle

\noindent\emph{Source and licence.} Compactness of Toeplitz Operators on the Bergman Space. Version arXiv:2609.24614v3 (CC BY 4.0), distributed by arXiv under the Creative Commons Attribution 4.0 International licence, http://creativecommons.org/licenses/by/4.0/, as recorded in the arXiv licence metadata for this exact version and shown on the abstract page https://arxiv.org/abs/2609.24614v3. Full licence notice: \texttt{licenses/arxiv-2609.24614-CC-BY-4.0.txt}.\par\smallskip

\noindent\textit{Editorial scope.} Counted unit: the article's lemma lem:halfplane-unique (source0.tex lines 559--580). The article's complete original proof of it is delivered with it (source0.tex lines 581--1020, one \texttt{proof} environment of 440 lines). No mathematics is written, repaired or re-derived: the statement and the proof are the article's own lines, in the article's own notation, and no retained line is substituted or re-flowed.  No further source-local context is needed: every cross-reference of every retained line resolves inside the two delivered spans.  Nothing was omitted from any retained span, so the package carries no omission record.  No retained line contains a citation, so the wrapper carries no bibliography and no bibliography entry.  No retained line contains a figure environment, an image-inclusion command or drawing code, so no image is delivered and the package declares no asset dependency.  Editorial formatting only, in addition: an AMS \texttt{amsart} preamble that restates the article's own declarations and macros used by the retained lines, namely the environments \texttt{lemma} on separate counters, together with the standard packages those lines need.  The retained environments print in this document as Lemma 1 in order of appearance, whereas the article prints the same items with the numbering of its own full document.  The complete native source of the article as distributed in the arXiv e-print for this exact version is bundled unchanged as \texttt{sources/211220/source0.tex}.
\begin{lemma}\label{lem:halfplane-unique}
Fix $0<q<1$. If $F\in L^\infty(\mathbb H)$ satisfies
\begin{equation}\label{eq:halfplane-zero}
\int_{B((x,t),qt)}F(u,v)\dd u\dd v=0
\qquad(x\in\mathbb R,\ t>0),
\end{equation}
then $F=0$ almost everywhere. Here
\[
B((x,t),qt)
=
\left\{(u,v)\in\mathbb{R}^2:
(u-x)^2+(v-t)^2<(qt)^2
\right\}
\]
denotes the Euclidean disk centered at $(x,t)$ with radius $qt$.
Since $0<q<1$, this disk is contained in the upper half-plane
$\mathbb H=\mathbb R\times(0,\infty)$. Consequently,
\begin{equation}\label{eq:halfplane-dense}
\overline{\Span\{\mathbf1_{B((x,t),qt)}:x\in\mathbb R,\ t>0\}}^{\,L^1(\mathbb H)}
=L^1(\mathbb H).
\end{equation}
\end{lemma}

\begin{proof}
We write $\mathcal S(\mathbb R)$ for the Schwartz space of
smooth functions $\vartheta$ such that
\[
\sup_{x\in\mathbb R}
(1+|x|)^m |\vartheta^{(n)}(x)|<\infty
\qquad (m,n\in\mathbb N_0),
\]
and $\mathcal S'(\mathbb R)$ for its continuous dual,
the space of tempered distributions. Choose a Schwartz function $\psi$ such that $\int\psi=1$ and $\widehat\psi\in C_c^\infty(\mathbb R)$.
Let $\psi_\varepsilon(x)=\varepsilon^{-1}\psi(x/\varepsilon)$, and define
\[
F_\varepsilon(x,t)=\int_{\mathbb R}\psi_\varepsilon(b)F(x-b,t)\dd b.
\]
Next, choose a nonnegative function $\eta_\delta\in C_c^\infty(-\delta,\delta)$ with $\int\eta_\delta=1$, and set
\begin{equation}\label{eq:single-regularize}
G(x,t)=\int_{\mathbb R}\eta_\delta(s)F_\varepsilon(e^sx,e^st)\dd s.
\end{equation}
For the limiting argument at the end of the proof, these functions are chosen from a single scaling family, namely
\[
\eta_\delta(s)=\delta^{-1}\eta(s/\delta),
\]
where $\eta\in C_c^\infty(-1,1)$ is nonnegative and
$\int_{\mathbb R}\eta(s)\dd s=1$. Thus $\{\eta_\delta\}_{\delta>0}$ is an
approximate identity on $\mathbb R$.
All integrals are justified by boundedness and the integrability of the kernels, and Fubini's theorem applies.
Horizontal translations preserve the family of disks, and dilations satisfy
$e^s B((x,t),qt)=B((e^sx,e^st),qe^st)$.
Hence both $F_\varepsilon$ and $G$ satisfy \eqref{eq:halfplane-zero}.
The dilation change of variables introduces only the nonzero factor $e^{-2s}$, which does not affect the vanishing of the integrals.

For completeness, the translation step follows from Fubini's theorem:
\begin{align*}
\int_{B((x,t),qt)}F_\varepsilon(u,v)\dd u\dd v
&=\int_{\mathbb R}\psi_\varepsilon(b)
  \int_{B((x,t),qt)}F(u-b,v)\dd u\dd v\dd b\\
&=\int_{\mathbb R}\psi_\varepsilon(b)
  \int_{B((x-b,t),qt)}F(y,v)\dd y\dd v\dd b=0.
\end{align*}
Likewise, using $(y,v)=(e^su,e^sv)$, we have
\begin{align*}
\int_{B((x,t),qt)}G(u,v)\dd u\dd v
&=\int_{\mathbb R}\eta_\delta(s)e^{-2s}
  \int_{B((e^sx,e^st),qe^st)}F_\varepsilon(y,v)\dd y\dd v\dd s=0.
\end{align*}
The absolute convergence needed for these applications of Fubini follows from
$F\in L^\infty(\mathbb H)$, $\psi_\varepsilon\in L^1(\mathbb R)$, and the
compact support of $\eta_\delta$.


Fix $\varepsilon,\delta>0$ from now on. Constants in the
regularity estimates below may depend on these parameters,
the chosen kernels, and $\norm F_\infty$, but not on $x$ or $t$.
Write
\[
H_j=\partial_x^jF_\varepsilon
     =(\partial_x^j\psi_\varepsilon)*_xF,
\qquad j\ge0.
\]
These horizontal derivatives are well defined distributionally,
with bounded measurable representatives satisfying
\[
\norm{H_j}_\infty
\le \norm F_\infty
   \norm{\partial_x^j\psi_\varepsilon}_{L^1(\mathbb R)}.
\]

We first verify that $G$ has a smooth representative.
Changing variables $v=e^st$ in
\eqref{eq:single-regularize} gives
\begin{equation}\label{eq:single-smooth-kernel}
G(x,t)
=
\int_0^\infty\int_{\mathbb R}
F(y,v)\,
\eta_\delta\!\left(\log\frac vt\right)
\psi_\varepsilon\!\left(\frac vt x-y\right)
\dd y\,\frac{\dd v}{v}.
\end{equation}
On each compact subset of $\mathbb H$, the $v$-integration
is restricted to a fixed compact subinterval of $(0,\infty)$.
Moreover, every derivative in $(x,t)$ of the kernel in
\eqref{eq:single-smooth-kernel} is dominated there by an
integrable function of $(y,v)$: the coefficients are bounded
on that compact set, and the horizontal kernels and all their
derivatives are Schwartz functions with uniformly bounded
translations. Since $F$ is bounded, differentiation under the
integral is valid to every order. Thus $G\in C^\infty(\mathbb H)$.

For the quantitative estimates, let
\[
D=x\partial_x+t\partial_t.
\]
Horizontal differentiation and integration by parts in the
scale variable, initially in the distributional sense, give
\begin{align*}
G_x
&=\int_{\mathbb R}\eta_\delta(s)e^s
       H_1(e^sx,e^st)\dd s,\\
G_{xx}
&=\int_{\mathbb R}\eta_\delta(s)e^{2s}
       H_2(e^sx,e^st)\dd s,\\
DG
&=-\int_{\mathbb R}\eta_\delta'(s)
       H_0(e^sx,e^st)\dd s,\\
D^2G
&=\int_{\mathbb R}\eta_\delta''(s)
       H_0(e^sx,e^st)\dd s,\\
D(G_x)
&=-\int_{\mathbb R}
       \bigl(\eta_\delta'(s)+\eta_\delta(s)\bigr)e^s
       H_1(e^sx,e^st)\dd s.
\end{align*}
Consequently,
\[
|G|+|G_x|+|G_{xx}|+|DG|+|D^2G|+|D(G_x)|
\le C.
\]
The identities
\[
tG_t=DG-xG_x,\qquad
tG_{xt}=D(G_x)-xG_{xx}
\]
and
\[
t^2G_{tt}
=D^2G-DG-2xD(G_x)+x^2G_{xx}
\]
therefore imply
\begin{equation}\label{eq:single-derivatives}
|G(x,t)|+|G_x(x,t)|\le C,
\qquad
|G_t(x,t)|\le \frac{C(1+|x|)}{t},
\end{equation}
as well as
\begin{equation}\label{eq:single-second-derivatives}
|G_{xx}(x,t)|\le C,\qquad
|G_{xt}(x,t)|\le \frac{C(1+|x|)}{t},\qquad
|G_{tt}(x,t)|\le \frac{C(1+x^2)}{t^2}.
\end{equation}
In particular, on every strip $0<a\le t\le b<\infty$,
all derivatives of total order at most two are bounded by
$C_{a,b}(1+x^2)$.

Finally, choose $L<\infty$ such that
$\operatorname{supp}\widehat{\psi_\varepsilon}\subset[-L,L]$.
Horizontal convolution gives this Fourier support bound for
$F_\varepsilon(\cdot,t)$ at almost every height.
Since $s\in\operatorname{supp}\eta_\delta\subset[-\delta,\delta]$,
horizontal dilation enlarges the support by at most $e^\delta$.
It follows that
\[
\operatorname{supp}\mathcal F_xG(\cdot,t)
\subset[-\Lambda,\Lambda],
\qquad \Lambda=e^\delta L,
\]
for every $t>0$. Here Fourier transforms are understood in
the sense of tempered distributions. Indeed, the uniform bound on $G$ and its pointwise continuity
imply that, for every $\vartheta\in\mathcal S(\mathbb R)$,
\[
t\longmapsto
\int_{\mathbb R}G(x,t)\vartheta(x)\dd x
\]
is continuous by dominated convergence. Hence
$G(\cdot,t)$ is continuous in the weak topology of
$\mathcal S'(\mathbb R)$. Since the Fourier transform is continuous
on $\mathcal S'(\mathbb R)$, the common Fourier support bound
extends from almost every height to every height by testing
$\mathcal F_xG(\cdot,t)$ against functions in
$C_c^\infty(\mathbb R\setminus[-\Lambda,\Lambda])$.

For fixed $\varepsilon$, let $\delta\downarrow0$. Continuity of dilations in local $L^1$ gives
$G\to F_\varepsilon$ in $L^1_{\mathrm{loc}}$. Then, as $\varepsilon\downarrow0$,
$F_\varepsilon\to F$ in $L^1_{\mathrm{loc}}$.
More explicitly, for every compact set $K\Subset\mathbb H$, the choice of
$\eta_\delta$ above and the continuity of dilations in $L^1(K)$ imply
\[
G\longrightarrow F_\varepsilon\quad\text{in }L^1(K)\qquad(\delta\downarrow0).
\]
The standard approximation-of-the-identity property of horizontal convolution
similarly gives
\[
F_\varepsilon\longrightarrow F\quad\text{in }L^1(K)\qquad(\varepsilon\downarrow0).
\]
Consequently, once the vanishing of $G$ has been established for every fixed
$\varepsilon,\delta>0$, these two local $L^1$ convergences imply $F=0$ almost everywhere.
Thus it suffices to prove that each $G$ vanishes.

\medskip
Consider the Hilbert space
\[
X=L^2(\mathbb R,w(x)\dd x),\qquad w(x)=(1+x^2)^{-4}.
\]
Then
\[
|\langle v,\vartheta\rangle|
\le\|v\|_X\left(\int|\vartheta(x)|^2(1+x^2)^4dx\right)^{1/2}, \quad v \in X,
\quad\vartheta\in\mathcal S,
\]
Bounded functions and functions of polynomial growth of degree at most two belong to $X$.
For $v\in X$ and $h\in\mathbb R$, define the horizontal
translation $\tau_hv$ by
\[
(\tau_hv)(x)=v(x-h).
\]
Since
\[
1+u^2\le 2\bigl(1+(u+h)^2\bigr)(1+h^2),
\]
we have
\[
w(u+h)\le 16(1+h^2)^4w(u).
\]
Consequently,
\[
\begin{aligned}
\|\tau_hv\|_X^2
&=\int_{\mathbb R}|v(u)|^2w(u+h)\,du\\
&\le 16(1+h^2)^4\|v\|_X^2,
\end{aligned}
\]
and hence
\[
\|\tau_hv\|_X
\le 4(1+h^2)^2\|v\|_X
\le 4(1+|h|)^4\|v\|_X.
\]
Choose $\chi\in C_c^\infty(\mathbb R)$ that is identically $1$ on a neighborhood of $[-\Lambda,\Lambda]$.
With the convention $\widehat v(\xi)=\int v(x)e^{-ix\xi}\dd x$, set
\[
K=\mathcal F^{-1}(-\xi^2\chi(\xi)),\qquad Bv=K*v.
\]
The function $K$ is Schwartz. Minkowski's inequality and the translation estimate above give
\[
\norm{Bv}_X\le
C\left(\int|K(h)|(1+|h|)^4\dd h\right)\norm v_X.
\]
Thus $B$ is bounded on $X$. If the horizontal Fourier support of $v$ is contained in
$[-\Lambda,\Lambda]$, then $Bv=\partial_x^2v$, initially in the distributional sense,
and also in $X$ for the smooth functions used below. 

\medskip
For $0\le r<t$, define
\[
U(x,t,r)=\frac1\pi\int_{a^2+b^2<1}G(x+ra,t+rb)\dd a\dd b.
\]
For each $(t,r)$, we regard $U(\cdot,t,r)$ as an element of $X$ and,
when no confusion can arise, write it simply as $U(t,r)$. All derivatives of
$U$ below are understood in this $X$-valued sense.
For $r>0$, this is the ordinary disk average of radius $r$,
and $U(x,t,0)=G(x,t)$.
For the regularity argument, use the same formula to extend
$U$ to the open parameter set
\[
\mathcal O=\{(t,r)\in\mathbb R^2:t>0,\ |r|<t\}.
\]
We claim that
\[
(t,r)\longmapsto U(\cdot,t,r)
\]
belongs to $C^2(\mathcal O;X)$, where derivatives and continuity
are taken in the norm of $X$.

Let $\mathcal K\Subset\mathcal O$. There exist constants
$a_{\mathcal K},b_{\mathcal K},L_{\mathcal K}>0$ such that,
for $(t,r)\in\mathcal K$ and $a^2+b^2\le1$,
\[
a_{\mathcal K}\le t+rb\le b_{\mathcal K},
\qquad |ra|\le L_{\mathcal K}.
\]
Thus \eqref{eq:single-derivatives} and
\eqref{eq:single-second-derivatives} show that all parameter
derivatives of the integrand of total order at most two
are bounded in absolute value by $C_{\mathcal K}(1+x^2)$.
Indeed, if
\[
\mathcal A H(x,t,r)
=\frac1\pi\int_{a^2+b^2<1}
H(x+ra,t+rb)\dd a\dd b,
\]
then the pointwise derivative formulas are
\begin{align*}
U_t&=\mathcal A(G_t),\\
U_r&=\frac1\pi\int_{a^2+b^2<1}
       (aG_x+bG_t)(x+ra,t+rb)\dd a\dd b,\\
U_{tt}&=\mathcal A(G_{tt}),\\
U_{tr}&=\frac1\pi\int_{a^2+b^2<1}
       (aG_{xt}+bG_{tt})(x+ra,t+rb)\dd a\dd b,\\
U_{rr}&=\frac1\pi\int_{a^2+b^2<1}
       (a^2G_{xx}+2abG_{xt}+b^2G_{tt})
       (x+ra,t+rb)\dd a\dd b.
\end{align*}

The dominating function belongs to $X$, since
\[
\int_{\mathbb R}(1+x^2)^2w(x)\dd x
=
\int_{\mathbb R}(1+x^2)^{-2}\dd x<\infty.
\]
The smoothness of $G$ and dominated convergence therefore
give continuity in $X$ of all the displayed expressions.
Applying the same domination to parameter difference
quotients, using the fundamental theorem of calculus on
a slightly larger compact subset of $\mathcal O$, proves
that these expressions are the strong first and second
derivatives of $U$. Hence $U\in C^2(\mathcal O;X)$.

The horizontal Fourier support of $U(\cdot,t,r)$ remains
contained in $[-\Lambda,\Lambda]$: horizontal translations
do not enlarge Fourier support, and the averaging only
combines functions having this common support.
Moreover,
\[
U_{xx}(x,t,r)
=\frac1\pi\int_{a^2+b^2<1}
G_{xx}(x+ra,t+rb)\dd a\dd b
\]
belongs to $X$. Consequently,
\[
BU(\cdot,t,r)=U_{xx}(\cdot,t,r)
\]
as an equality in $X$. Applying \eqref{eq:halfplane-zero} to $G$ at the center $(x,t)$ and radius
$qt$ gives, for every $t>0$,
\begin{equation}\label{eq:single-movingboundary}
U(t,qt)=0.
\end{equation}
Disk averages satisfy
\begin{equation}\label{eq:single-darboux}
U_{rr}+\frac3rU_r=U_{xx}+U_{tt}.
\end{equation}
Let $S(x,t,r)$ be the circular average of $G$. The divergence theorem gives
$U_{xx}+U_{tt}=2S_r/r$.
Since $U=2r^{-2}\int_0^r sS(x,t,s)\dd s$, we obtain
$S=U+(r/2)U_r$, which yields \eqref{eq:single-darboux} upon substitution.
Consequently, in $X$,
\begin{equation}\label{eq:single-wave}
U_{tt}=U_{rr}+\frac3rU_r-BU.
\end{equation}
The identity $U_r(t,0)=0$ follows from the vanishing of the first moments of the unit disk.

\medskip
The weight $r^3$ (i.e.\ $r^{n+1}$ with $n=2$) is chosen so that the radial part $U_{rr}+\tfrac3rU_r$ is  symmetric on $L^2(r^3\,\dd r)$,
which eliminates the boundary terms in the integration by parts below.
Define the nonnegative energy
\begin{equation}\label{eq:single-energy}
E(t)=\frac12\int_0^{qt}
\bigl(\norm{U_t(t,r)}_X^2+\norm{U_r(t,r)}_X^2+\norm{U(t,r)}_X^2\bigr)r^3\dd r.
\end{equation}
On every compact interval $0<a\le t\le b$, the preceding regularity justifies differentiation of the energy and integration by parts.
To make the behavior at the lower endpoint explicit, for $0<\alpha<qt$ first
set
\[
E_\alpha(t)=\frac12\int_\alpha^{qt}
\bigl(\norm{U_t(t,r)}_X^2+\norm{U_r(t,r)}_X^2+\norm{U(t,r)}_X^2\bigr)r^3\dd r.
\]
Differentiate $E_\alpha$ and integrate by parts in $r$. The terms at
$r=\alpha$ are bounded by a constant times $\alpha^3$ on $[a,b]$ and hence
vanish as $\alpha\downarrow0$.
The same $C^2(\mathcal O;X)$ regularity implies that $U_t$ and
$U_r$ remain bounded in $X$ as $r\downarrow0$, uniformly for
$t\in[a,b]$. Hence
\[
r^3\operatorname{Re}\langle U_t(t,r),U_r(t,r)\rangle_X
\longrightarrow0
\qquad(r\downarrow0),
\]
uniformly for $t\in[a,b]$. Let $\alpha\downarrow0$.
By \eqref{eq:single-wave},
\begin{align*}
E'(t)={}&(qt)^3\operatorname{Re}\langle U_t(t,qt),U_r(t,qt)\rangle_X\\
&+\frac q2(qt)^3\bigl(\norm{U_t(t,qt)}_X^2+
\norm{U_r(t,qt)}_X^2+\norm{U(t,qt)}_X^2\bigr)\\
&+\int_0^{qt}\operatorname{Re}\langle U_t,(I-B)U\rangle_Xr^3\dd r.
\end{align*}
Differentiating along the boundary in \eqref{eq:single-movingboundary} gives
$U_t(t,qt)+qU_r(t,qt)=0$.
Hence the first two lines equal
\[
-\frac{q(1-q^2)}2(qt)^3\norm{U_r(t,qt)}_X^2\le0.
\]
For the last term, the boundedness of $B$ on $X$ and Cauchy--Schwarz give
\begin{align*}
\left|\operatorname{Re}\langle U_t,(I-B)U\rangle_X\right|
&\le (1+\norm B)\norm{U_t}_X\norm U_X\\
&\le \frac{1+\norm B}{2}
   \bigl(\norm{U_t}_X^2+\norm U_X^2\bigr),
\end{align*}
where we used $2ab\leq a^2+b^2$. Since the integral of
$\norm{U_t}_X^2+\norm U_X^2$ against $r^3\dd r$ is bounded by
$2E(t)$, we may set
\[
C_0:=1+\norm{B}_{\mathcal B(X)}.
\]
Thus
\begin{equation}\label{eq:single-gronwall}
E'(t)\leq C_0 E(t),\qquad t>0.
\end{equation}
For fixed regularization parameters, $C_0$ is independent of $t$ and of the
endpoints of the interval on which the energy identity is justified.
\medskip

For $0<t\le1$ and $0\le r\le qt$, the sampling points satisfy
$t+rb\ge(1-q)t$ and $|x+ra|\le|x|+1$. By
\eqref{eq:single-derivatives},
\[
\norm{U(t,r)}_X\le C,
\qquad
\norm{U_t(t,r)}_X+\norm{U_r(t,r)}_X\le C/t.
\]
Here $U_r$ is the average of $aG_x+bG_t$, and the functions of at most
linear growth in $x$ that occur here belong to $X$. Therefore
\[
0\le E(t)\le C\int_0^{qt}(1+t^{-2})r^3\dd r
\le C(t^4+t^2)\longrightarrow0
\qquad(t\downarrow0).
\]
For fixed $T>0$ and every $0<s<T$, \eqref{eq:single-gronwall} gives
\[
E(T)\le e^{C_0(T-s)}E(s).
\]
Letting $s\downarrow0$ and using $E(s)\to0$ yields $E(T)=0$. Since the energy contains the term $\norm U_X^2$, continuity gives $U(t,r)=0$ for $0<r<qt$.
More precisely, $U(\cdot,T,r)$ is continuous as an $X$-valued function of
$r$, and $E(T)=0$ implies $U(\cdot,T,r)=0$ for every $0\le r\le qT$.

Letting $r\downarrow0$, we obtain $G(\cdot,t)=0$ in $X$.
Since $G$ is smooth, $G=0$ on $\mathbb H$.
Letting first $\delta\downarrow0$ and then $\varepsilon\downarrow0$, we obtain $F=0$ almost everywhere.
\medskip

If \eqref{eq:halfplane-dense} were false, the Hahn--Banach theorem would
produce a nonzero continuous linear functional on $L^1(\mathbb H)$ that
annihilates every disk indicator. Since
$L^1(\mathbb H)^*=L^\infty(\mathbb H)$, this functional is represented by a
nonzero $F\in L^\infty(\mathbb H)$. Hence
\[
\int_{B((x,t),qt)}F(u,v)\dd u\dd v=0
\qquad(x\in\mathbb R,\ t>0),
\]
which contradicts the uniqueness just proved.
\end{proof}

\end{document}
